% Generated draft: requires mathematical and source review.
\documentclass[11pt]{article}
\usepackage[margin=25mm]{geometry}
\usepackage{fontspec}
\setmainfont{Latin Modern Roman}
\usepackage{amsmath,amssymb,mathtools,unicode-math}
\setmathfont{Latin Modern Math}
\usepackage{enumitem,needspace,xurl,hyperref}
\setlength{\parindent}{0pt}
\setlength{\parskip}{5pt}
\setlength{\emergencystretch}{3em}
\newcommand{\N}{\mathbb{N}}
\newcommand{\Z}{\mathbb{Z}}
\newcommand{\Q}{\mathbb{Q}}
\newcommand{\R}{\mathbb{R}}
\newcommand{\C}{\mathbb{C}}
\newcommand{\F}{\mathbb{F}}
\newcommand{\A}{\mathbb{A}}
\newcommand{\PP}{\mathbb{P}}
\newcommand{\E}{\mathbb{E}}
\newcommand{\eps}{\varepsilon}
\newcommand{\dd}{\,\mathrm{d}}
\begin{document}
\section*{UnsolvedMath ID: 30001089}
\textbf{Source problem code:} catalogue entry
\subsection{Definitions and mathematical statement}
This entry is a bundle of questions in complex dimension two. A rotation domain is understood in the report's dynamical setting, with a suitable iterate analytically conjugate to a rotation; a Siegel ball is the corresponding ball-type domain. The ambient map class, rotation rank, iterate, and exact ball convention are not reproduced. Retain the following distinct targets:
\begin{enumerate}
\item Does a Siegel ball exist whose boundary is homeomorphic to \(S^3\)? The requested boundary property is topological, with no smoothness asserted.
\item Does a rotation domain exist that is biholomorphic to \(A\times\C\), where \(A\) is an annulus? The annulus modulus and allowed map class are unspecified.
\item Describe the domain of convergence of the linearizing map for a two-dimensional rotation domain. Neither the chosen linearization nor a proposed geometric description is given.
\item For a complex H\'enon map with connected Julia set \(J\), when is the report's natural map \(J^-\setminus K\to\Sigma\) to the complex solenoid one-to-one? When is \(J\) a quotient of a solenoid, and which equivalence relations on that solenoid occur?
\item Find an effective structural characterization of connectedness of the Julia set for H\'enon-like maps; the source gives no specific proposed criterion.
\end{enumerate}
For orientation, for an invertible map \(f:\C^2\to\C^2\) of H\'enon type, customary notation is
\[
K^\pm=\{z:\{f^{\pm n}(z):n\geq0\}\text{ is bounded}\},\qquad
J^\pm=\partial K^\pm,\qquad
K=K^+\cap K^-,\quad J=J^+\cap J^-.
\]
The report's conventions must be confirmed. In particular, the construction and domain of the ``natural map,'' the precise complex solenoid \(\Sigma\), its degree, and the allowed quotient equivalence relations are absent. A complex solenoid is not silently replaced by a compact real circle solenoid, and connectedness alone is not supplemented by hyperbolicity. The H\'enon-like class in the last question is likewise not defined here. These omissions prevent turning the bundle into one self-contained theorem or assigning one verified status to all its parts.
\subsection{Short English statement}
Ask about sphere-boundary Siegel balls, annulus-times-plane rotation domains, linearization domains, solenoid descriptions for connected H\'enon Julia sets, and connectedness criteria for H\'enon-like maps. Several map and solenoid definitions are missing.
\subsection{Sources}
\begin{itemize}
\item Local UnsolvedMath registry, permanent ID \(30001089\), source code catalogue entry, ``Rotation Domains and Solenoid Models in Higher-Dimensional Dynamics.'' No individual proposer is supplied.
\item \emph{Trends and Developments in Complex Dynamics} (2009), Oberwolfach Reports, DOI 10.4171/owr/2008/48, pp. 2703--2762: \url{https://doi.org/10.4171/owr/2008/48}. This is the sole supplied external reference.
\item Supplied literature triage, checked 2026-08-21: status ``open.'' It describes several proposed problems rather than one complete claim and reports no authoritative source establishing all of them. This is supplied dated curation, not independently verified; it does not individually certify every subproblem open, and time-sensitive status and cited claims require independent verification.
\end{itemize}
\end{document}
